\documentclass[withoutpreface,bwprint]{cumcmthesis}

\title{大型展销会临时工招聘与排班优化}
\tihao{【待填写题号】}
\baominghao{【待填写报名号】}
\schoolname{【待填写学校名称】}
\membera{【待填写队员一】}
\memberb{【待填写队员二】}
\memberc{【待填写队员三】}
\supervisor{【待填写指导教师】}
\yearinput{2026}
\monthinput{【待填写】}
\dayinput{【待填写】}

\newcolumntype{Y}{>{\centering\arraybackslash}X}
\newcolumntype{L}{>{\raggedright\arraybackslash}X}
\newcommand{\R}{\mathbb{R}}
\newcommand{\Z}{\mathbb{Z}}
\newcommand{\ind}{\mathbf{1}}

% 显式待补充标记：正式定稿时逐项删除或替换为正文。
\newcommand{\pending}[1]{%
  \par\medskip
  \noindent\fbox{%
    \parbox{0.94\linewidth}{\textbf{【待补充】}#1}%
  }%
  \par\medskip
}

% 保留原图文件名、尺寸、题注和版面位置。
% 原图未上传时仅显示占位框；将同名图片放入 figures/ 后自动恢复。
\newcommand{\originalfigure}[2][]{%
  \IfFileExists{figures/#2}{\includegraphics[#1]{figures/#2}}{%
  \IfFileExists{figures/#2.png}{\includegraphics[#1]{figures/#2.png}}{%
  \IfFileExists{figures/#2.pdf}{\includegraphics[#1]{figures/#2.pdf}}{%
  \IfFileExists{figures/#2.jpg}{\includegraphics[#1]{figures/#2.jpg}}{%
    \fbox{\parbox[c][4.2cm][c]{0.88\linewidth}{\centering
      \textbf{原图暂未提供}\par
      请将 \texttt{\detokenize{#2}} 放入 \texttt{figures/}\par
      图片内容、尺寸与题注暂不调整。}}
  }}}}
}

\begin{document}

\maketitle

\begin{abstract}
针对为期10天、包含10个工作小组和11个营业小时的大型展销会临时工招聘与排班问题，本文在不改变题目三种工作模式的前提下，设计“宏观班型优化--微观员工指派”的双层求解框架。宏观层预先枚举合法匿名班型，以混合整数线性规划确定每天各小组使用各类班型的人数；微观层先用CP-SAT确定每名员工的工作与休息状态，再将宏观班型确定性映射到具体员工。该分解避免直接对“员工--日期--小组--班型”进行高维联合枚举，同时保留最优性和可执行性。

\textbf{对于问题一}，在各小组人员固定的条件下，枚举三种工作模式对应的合法班型，并对每个“日期—小组”建立逐小时集合覆盖整数规划，求得日内最低工作人数。根据单日峰值需求与十天总工作人日约束，确定十个小组的最少人员规模分别为 39，39，43，41，44，39，42，42，43，45，共需招聘 \textbf{417人}。进一步利用最大流安排员工休息日，并将每日在岗员工与班次一一对应，从而实现每人工作8天、休息2天的可执行排班。

\textbf{对于问题二}，将十个固定人员池合并为统一人员池，汇总得到逐日最低工作人数总和为3247。综合最高单日需求和每名员工必须工作8天的约束，确定最少招聘人数为 \textbf{406人}，较问题一减少11人。由于406名员工必须提供 $406 \times 8 = 3248$ 个工作人日，实际排班需在逐日覆盖下界基础上增加1个整除冗余人日，模型将其安排在第7天。随后采用最大流构造员工工作—休息矩阵，并完成员工与具体班次的确定性映射，使每名员工均恰好工作8天、休息2天，且所有日期、小组和时段的人员需求均得到满足。

\textbf{对于问题三}，本文不把求解器返回的“不可行”当作终点，而是给出公共盲区定理。营业窗口长度为11小时，两个4小时工作段之间至少休息2小时，因此全部合法班型只有三种，且共同无法覆盖13:00--14:00；附件在该时段十天合计需求为1500人次，故原问题具有与招聘人数无关的结构性不可行性。对此比较三类规则修复：将休息缩短至1小时需要437名全职员工；采用非对称4+3班型需495人，同时允许3+4班型可降至439人；主推的混合用工方案保留全职员工“4小时工作+至少2小时休息+4小时工作”，只以4小时兼职班修补午间盲区。采用“先最小化兼职人次、再最小化全职人数”的字典序目标，得到兼职理论下界\textbf{1500人次}和全职最优规模\textbf{325人}。按同小时工资折算，等效成本为418.75名全职员工，付费工时浪费率为34.02\%，在满足严格休息制度的可行方案中同时改善成本与人效。

模型由SCIP求得整数最优解，并以CP-SAT、逐小时覆盖复核、员工级人日对账、工资权重敏感性和多方案二维评价进行验证。结果说明，排班规则的结构性分析比盲目增加人数更重要，而有限、定向的兼职柔性可以打破严格全职制度造成的可行性与成本锁死。

\pending{待补充：摘要末尾尚需用 1--2 句话独立凝练“公共盲区定理、宏观--微观双层排班与字典序混合用工”等创新点。}

\keywords{人员排班；混合整数规划；CP-SAT；结构性不可行；混合用工；字典序优化}
\end{abstract}

\section{问题重述}

\subsection{问题背景}

大型展销会通常具有举办周期短、服务环节多和客流波动明显等特点。为保障咨询引导、秩序维护及现场服务等工作正常开展，主办方需要根据不同日期、不同时间段及不同工作区域的用工需求，提前完成临时人员的招聘与排班。若招聘人数过少，可能造成部分岗位缺员，影响展销会的服务质量与运行秩序；若招聘人数过多，则会增加人工成本，并产生较多闲置工时。因此，如何在满足各时段最低用工需求和员工劳动规则的前提下，合理确定招聘规模并制定可执行的排班方案，具有较强的实际管理价值。

人员排班问题本质上是在时变需求、劳动制度和用工成本之间进行协调的组合优化问题，已广泛应用于零售、交通、医疗和公共服务等领域\cite{ernst2004,vandenbergh2013personnel}。经典班次排班研究通常将员工可采用的工作时间组合表示为若干合法班型，并通过确定各班型的配置人数覆盖不同时段的人员需求\cite{dantzig1954comment}。然而，当排班周期跨越多个日期，且员工还需满足规定的工作天数、休息天数和岗位归属要求时，日内班次安排与跨日人员配置会相互耦合，使问题的求解难度显著增加。

本题研究一个为期10天的大型展销会临时工招聘与排班问题。展销会设置10个工作小组，每日营业时间为08:00--19:00，附件给出了各日期、各小时和各小组的最低用工需求，共形成1100个需求单元。每名员工在会期内工作8天、休息2天，并在每个工作日完成两个连续4小时的工作段。问题一要求员工在整个会期内固定服务于同一小组，问题二则允许员工在不同日期更换服务小组。两种人员组织方式分别对应固定人员池和统一人员池，其核心均是协调逐小时需求覆盖与跨日人员共享，在保证排班可执行的基础上求得最少招聘人数。


\subsection{问题要求}

根据附件给出的10天逐小时用工需求，题目要求在满足临时工工作时间和休息规则的条件下，研究不同人员组织方式下的招聘与排班问题。具体要求如下。

\textbf{问题一}要求整个展销会期间固定服务于同一个工作小组，不得在不同小组之间调动。每名员工在10天内工作8天、休息2天，并在每个工作日完成两个连续4小时的工作段。要求根据各日期、各小时和各小组的最低用工需求，分别确定10个小组的招聘人数，制定具体排班方案，并使全场招聘总人数最少。

\textbf{问题二}要求临时工在同一天内只能服务于一个工作小组，但允许在不同日期被分配至不同小组。每名员工在10天内工作8天、休息2天，并在每个工作日完成两个连续4小时的工作段。要求在统一调配人员的条件下确定全场最少招聘人数，并制定具体排班方案，明确每名员工的工作日期、服务小组和工作时段，使各日期、各小时及各小组的最低用工需求均得到满足。

\textbf{问题三}要求临时工在同一天内可以服务至多两个工作小组，每个工作段连续工作4小时，且两个工作段之间至少休息2小时。要求判断该工作制度下是否存在满足全部逐小时用工需求的排班方案；若原有规则无法形成可行排班，则分析不可行的原因，并提出合理的用工制度调整方案，在满足需求覆盖和员工休息要求的基础上，对不同方案的人员规模、用工成本及人员利用效率进行比较。


\section{模型假设}

为明确模型的适用范围并简化实际排班过程，作如下假设：

\begin{enumerate}[label=\textbf{A\arabic*:}, wide=\parindent]

  \item \textbf{需求准确性：}假设附件给出的各日期、各小时及各小组的用工人数准确可靠，并将其视为对应时段必须满足的最低在岗人数，不考虑客流变化及其他突发因素引起的临时需求波动。

  \item \textbf{时段离散性：}假设每个小时内的人员需求保持不变，所有工作段均在整点开始或结束。员工在某一小时内完整在岗时，为对应小组提供一个单位的人员覆盖。

  \item \textbf{员工同质性：}假设各临时工的工作能力和服务效率相同，均能胜任所分配小组的工作，不考虑员工经验、熟练程度、技能等级及个体效率差异。

  \item \textbf{排班执行一致性：}假设所有员工均严格按照排班方案到岗，不发生迟到、早退、请假、中途离职或临时换班等情况，因而计划排班人数与实际在岗人数一致。

  \item \textbf{问题一、二工作规则：}问题一和问题二中，每名员工在一个工作日内完成两个连续4小时的工作段。由于题目未规定两个工作段之间的最短休息时间，允许后一工作段紧接前一工作段。

  \item \textbf{问题三工作规则：}问题三中，每名全职员工在一个工作日内仍完成两个连续4小时的工作段，两个工作段之间至少休息2小时；同一工作段只能服务于一个小组，但前后两个工作段可以分别服务于不同小组。

  \item \textbf{全职员工工作周期：}每名全职员工在10天会期内均恰好工作8天、休息2天；每个工作日只能接受一个由两个工作段组成的完整班型，不考虑额外加班和不足8小时的全职工作日。

  \item \textbf{兼职人员定位：}在问题三的混合用工方案中，兼职人员每次连续工作4小时，同一兼职班次只服务于一个工作小组。兼职人员用于补充严格全职排班难以覆盖的局部需求，不作为全面替代全职员工的常规人员池。

  \item \textbf{成本口径：}在不同方案的成本比较中，假设全职员工与兼职员工的单位小时工资相同。一名全职员工在会期内提供
  \[
  8\times8=64
  \]
  个付费工时，一个兼职班提供4个付费工时；暂不考虑招聘固定成本、培训成本、跨组调配成本、工资溢价及其他管理费用。

  \item \textbf{附加约束简化：}除题目明确规定的工作天数、休息天数、工作段长度和小组调配范围外，不考虑员工个人偏好、指定休息日期、连续工作天数上限及排班公平性等附加限制。

\end{enumerate}

\section{符号说明}

本文将08:00--09:00记为小时索引$h=0$，依次至
18:00--19:00的$h=10$。其中，$d,h,g,k,i,s,t$分别表示
日期、小时、小组、员工、班型及兼职班起始时段的索引，
$g$和$k$均表示工作小组。主要符号见表\ref{tab:symbols}。

\begin{table}[H]
\centering
\caption{主要集合、参数与变量}
\label{tab:symbols}

\renewcommand{\arraystretch}{1.15}

\begin{tabularx}{\textwidth}{
  >{\centering\arraybackslash}p{2.1cm}
  >{\centering\arraybackslash}X
  >{\centering\arraybackslash}p{3.0cm}
}
\toprule
\textbf{符号} & \textbf{含义} & \textbf{取值或单位}\\
\midrule

$D$
& 日期集合
& $\{1,2,\ldots,10\}$\\

$H$
& 小时索引集合
& $\{0,1,\ldots,10\}$\\

$G$
& 工作小组集合
& $\{1,2,\ldots,10\}$\\

$q_{dhg}$
& 第$d$天小时$h$小组$g$的最低需求
& 人\\

$S_r$
& 两段4小时工作之间至少间隔$r$小时的合法班型集合
& 集合\\

$c_{hs}$
& 班型$s$是否覆盖小时$h$
& $0$或$1$\\

$W_g,W$
& 问题一中小组$g$的固定编制、问题二中的统一编制
& 非负整数（人）\\

$x_{dgs}$
& 第$d$天小组$g$采用班型$s$的人数
& 非负整数（人）\\

$R_{dg},R_d$
& 问题一中各小组的休息人数、问题二中全场的休息人数
& 非负整数（人）\\

$y_{id}$
& 员工$i$在第$d$天是否工作
& $0$或$1$\\

$z_{dsgk}$
& 问题三中第$d$天采用班型$s$，前段服务小组$g$、
后段服务小组$k$的人数
& 非负整数（人）\\

$p_{dtg}$
& 第$d$天小组$g$使用起始时段$t$的4小时兼职人数
& 非负整数（人）\\

$P,W^F$
& 十天兼职班总人次、全职员工招聘人数
& 人次、人\\

\bottomrule
\end{tabularx}
\end{table}

\section{问题分析}

\subsection{问题一分析}

问题一要求在临时工\textbf{固定归组}的条件下，确定各小组的最少招聘人数，并制定满足逐小时需求和“做8休2”规则的排班方案。由于员工在整个会期内不能跨组调配，10个工作小组分别形成独立的人员池，因此可将原问题分解为\textbf{10个相互独立的排班子问题}。

首先，根据每名员工每天完成两个连续4小时工作段的规定，枚举营业时间内的全部\textbf{合法班型}，并建立各班型对11个小时的覆盖关系。对每个日期和工作小组分别建立\textbf{日内整数覆盖模型}，在满足各小时最低需求的条件下，求得当日所需的最少工作人数。利用合法班型覆盖分时需求是班次排班问题中的经典建模方法
\cite{dantzig1954comment,alfares2004tour}。

随后，在各组十天日内排班结果的基础上，同时考虑\textbf{单日工作人数峰值和十天总工作人日}，确定各小组\textbf{固定人员池的最小规模}。该阶段解决匿名人员总量问题，但所得总工作人日不能直接保证每名员工均恰好工作8天、休息2天。

因此，进一步对每个小组建立员工与日期之间的\textbf{最大流网络}，将每名员工的两个休息日分配到两个不同日期，并使每天的休息人数与日内排班结果一致
\cite{ford1956maximal}。在最大流达到全部休息人日后，再将每日在岗员工与对应的匿名班次槽位逐一映射，从而形成\textbf{员工级排班方案}。最后，利用完整混合整数规划对各组最少招聘人数进行独立复核。

\subsection{问题二分析}

问题二允许临时工在不同日期服务于不同工作小组，但同一天内仍只能服务于一个小组。与问题一相比，原有的10个固定人员池可以合并为统一人员池，使不同小组在不同日期出现的需求峰值能够通过跨日换组相互补偿，从而减少固定归组造成的人员冗余。

由于员工在同一天内仍只执行一个小组的完整班型，各日期、各小组的日内需求覆盖方式不发生改变。因此，首先沿用问题一的日内整数覆盖模型，分别求得各日期、各小组所需的最低工作人数，再将其汇总为每天全场的最低工作人数。

随后，根据十天中的最高单日工作人数和全部工作人日需求确定统一人员池规模。由于每名员工必须恰好工作8天，最终排班的总工作人日必须为招聘人数的8倍，而逐日独立求得的最低工作人日总和不一定满足这一整除条件。为此，需要在逐日覆盖下界的基础上增加最少的冗余工作人日，并确定其分配日期。

在得到各日实际工作人数后，建立统一人员池的员工--日期最大流网络，为每名员工分配两个不同的休息日，并保证各日休息人数与宏观方案一致。最后，将每天的在岗员工与各小组的匿名班次槽位进行确定性映射，使每名员工每天只获得一个完整班型，并自然满足同一天只服务一个小组的要求。

通过比较固定人员池与统一人员池的最少招聘人数，可进一步量化跨日人员共享对降低招聘规模和提高人员利用率的作用
\cite{hopp2004skill,taskiran2017cross}。

\subsection{问题三分析}


\section{问题一模型的建立与求解}

\subsection{合法班型的构造}

问题一和问题二均规定，每名临时工在一个工作日内完成两个连续4小时的工作段。为便于描述，将08:00--09:00记为小时索引$h=0$，依次至18:00--19:00的$h=10$，从而将每日营业时间划分为11个连续时段。

设班型$s=(u_s,v_s)$，其中$u_s$和$v_s$分别表示前、后两个工作段的开始索引。问题一和问题二未规定两个工作段之间的最短休息时间，因此两个工作段可以首尾相接，但不能发生重叠。合法班型集合为
\begin{equation}
S_0=
\left\{
(u,v)\in\mathbb{Z}^2:
0\leq u\leq3,\quad
u+4\leq v\leq7
\right\}.
\label{eq:q12-shift-set}
\end{equation}

由式\eqref{eq:q12-shift-set}共得到10种合法班型，具体见表\ref{tab:q12-shifts}。

\begin{table}[H]
\centering
\caption{问题一和问题二的合法班型}
\label{tab:q12-shifts}
\begin{tabular}{cccccc}
\toprule
班型 & 第一工作段 & 第二工作段
& 班型 & 第一工作段 & 第二工作段\\
\midrule
$s_1$ & 08:00--12:00 & 12:00--16:00
& $s_6$ & 09:00--13:00 & 14:00--18:00\\

$s_2$ & 08:00--12:00 & 13:00--17:00
& $s_7$ & 09:00--13:00 & 15:00--19:00\\

$s_3$ & 08:00--12:00 & 14:00--18:00
& $s_8$ & 10:00--14:00 & 14:00--18:00\\

$s_4$ & 08:00--12:00 & 15:00--19:00
& $s_9$ & 10:00--14:00 & 15:00--19:00\\

$s_5$ & 09:00--13:00 & 13:00--17:00
& $s_{10}$ & 11:00--15:00 & 15:00--19:00\\
\bottomrule
\end{tabular}
\end{table}

定义班型$s$对小时$h$的覆盖参数为
\begin{equation}
c_{hs}
=
\ind\{u_s\leq h<u_s+4\}
+
\ind\{v_s\leq h<v_s+4\}.
\label{eq:q12-shift-cover}
\end{equation}

由于两个工作段互不重叠，$c_{hs}$仅取0或1。由此可构造大小为$11\times10$的班型覆盖矩阵，并将逐小时需求覆盖问题转化为整数规划问题。

\subsection{日内最低用工模型}

对于给定的日期$d$和小组$g$，以$x_s$表示当天采用班型$s$的临时工人数。为求得满足该小组11个时段需求所需的最少工作人数，建立日内整数覆盖模型：
\begin{equation}
\begin{aligned}
m_{dg}=\min\quad
&\sum_{s\in S_0}x_s\\
\mathrm{s.t.}\quad
&\sum_{s\in S_0}c_{hs}x_s
\geq q_{dhg},
&&h\in H,\\
&x_s\in\mathbb{Z}_{\geq0},
&&s\in S_0.
\end{aligned}
\label{eq:q1-daily-cover}
\end{equation}

目标函数表示第$d$天小组$g$的工作人数，需求覆盖约束保证每个时段的在岗人数均不低于附件给出的最低需求。分别对10个日期和10个小组求解式\eqref{eq:q1-daily-cover}，可得到100个日内最低工作人数$m_{dg}$及其对应的最优班型组合。

每个子问题仅包含10个非负整数变量和11个需求覆盖约束，因此可以相互独立地求解。本文通过PySCIPOpt调用SCIP完成整数规划求解
\cite{achterberg2009scip}。

\subsection{固定人员规模的确定}

问题一规定员工在整个会期内固定服务于同一小组。对于第$g$组，固定人员规模首先不能低于任一日期的最低工作人数，即
\begin{equation}
W_g\geq\max_{d\in D}m_{dg}.
\label{eq:q1-peak-bound}
\end{equation}

另一方面，每名员工在10天内恰好工作8天，因此$W_g$名员工共能提供$8W_g$个工作人日。为覆盖该组十天的日内最低工作人日，必须满足
\begin{equation}
8W_g\geq\sum_{d\in D}m_{dg}.
\label{eq:q1-workday-bound}
\end{equation}

综合式\eqref{eq:q1-peak-bound}和式\eqref{eq:q1-workday-bound}，得到第$g$组固定招聘人数：
\begin{equation}
W_g^*
=
\max
\left\{
\max_{d\in D}m_{dg},
\left\lceil
\frac{\displaystyle\sum_{d\in D}m_{dg}}{8}
\right\rceil
\right\}.
\label{eq:q1-staff}
\end{equation}

由于每名员工必须恰好工作8天，第$g$组最终需要安排的工作人日为$8W_g^*$。日内独立最优结果与该数值之间的差额为
\begin{equation}
E_g
=
8W_g^*-\sum_{d\in D}m_{dg}.
\label{eq:q1-extra-workdays}
\end{equation}

设第$g$组第$d$天的实际工作人数为$A_{dg}$，则需要满足
\begin{align}
m_{dg}\leq A_{dg}&\leq W_g^*,
&&d\in D,\label{eq:q1-actual-bound}\\
\sum_{d\in D}A_{dg}&=8W_g^*.
\label{eq:q1-actual-total}
\end{align}

各日期可以容纳的补充工作人日总数为
\begin{equation}
\sum_{d\in D}\left(W_g^*-m_{dg}\right)
=
10W_g^*-\sum_{d\in D}m_{dg},
\end{equation}
且有
\[
10W_g^*-\sum_{d\in D}m_{dg}
\geq
8W_g^*-\sum_{d\in D}m_{dg}=E_g.
\]
因此，所需补充工作人日一定能够在不超过每日固定编制的条件下完成分配，式\eqref{eq:q1-staff}给出的招聘规模可以达到。

本文按照日内最低人数由大到小的顺序分配补充工作人日；当最低人数相同时，优先选择日期编号较小的日期。确定$A_{dg}$后，重新求解
\begin{align}
\sum_{s\in S_0}c_{hs}x_{dgs}
&\geq q_{dhg},
&&h\in H,\label{eq:q1-final-cover}\\
\sum_{s\in S_0}x_{dgs}
&=A_{dg},\label{eq:q1-final-count}\\
x_{dgs}&\in\mathbb{Z}_{\geq0},
&&s\in S_0,
\end{align}
并按班型编号进行确定性择优，得到最终的匿名班型配额。

\subsection{休息日期的最大流分配}

上述模型确定了各组每天的工作人数和班型配额，但尚未说明每名员工在哪两天休息。第$g$组第$d$天的休息人数为
\begin{equation}
R_{dg}=W_g^*-A_{dg}.
\label{eq:q1-rest-number}
\end{equation}

由式\eqref{eq:q1-actual-total}可得
\begin{equation}
\sum_{d\in D}R_{dg}
=
10W_g^*-\sum_{d\in D}A_{dg}
=
2W_g^*.
\label{eq:q1-rest-total}
\end{equation}

为将休息人日分配到具体员工，对每个小组分别建立“源点--员工节点--日期节点--汇点”的最大流网络。其边容量设置为
\begin{equation}
\begin{aligned}
\operatorname{cap}(o,e_i)&=2,\\
\operatorname{cap}(e_i,d)&=1,\\
\operatorname{cap}(d,t)&=R_{dg},
\end{aligned}
\label{eq:q1-flow-capacity}
\end{equation}
其中$o$和$t$分别表示源点和汇点，$e_i$表示组内员工节点。

源点到员工节点的容量2表示每名员工需要安排两个休息日；员工到日期节点的容量1保证同一员工在同一天至多休息一次；日期节点到汇点的容量$R_{dg}$保证第$d$天恰有$R_{dg}$名员工休息。

若最大流达到
\begin{equation}
F_g=2W_g^*,
\label{eq:q1-full-flow}
\end{equation}
则每名员工均获得两个不同的休息日，且各日期的休息人数与匿名排班方案完全一致。由于网络中全部容量均为整数，根据最大流的整数性，员工与日期之间的流量可直接转化为0--1休息状态。其余8天即为该员工的工作日。

\subsection{员工--班次确定性映射}

最大流确定工作日期后，记第$d$天第$g$组的在岗员工集合为
\[
\mathcal E_{dg},
\qquad
|\mathcal E_{dg}|=A_{dg}.
\]

将当天的匿名班型配额展开为班次槽位集合
\begin{equation}
\mathcal L_{dg}
=
\bigcup_{s\in S_0}
\underbrace{\{s,s,\ldots,s\}}_{x_{dgs}\text{个}},
\label{eq:q1-slot-set}
\end{equation}
则
\[
|\mathcal L_{dg}|
=
\sum_{s\in S_0}x_{dgs}
=
A_{dg}
=
|\mathcal E_{dg}|.
\]

因此，可以将当天在岗员工与匿名班次槽位进行一一对应。本文按照员工编号和班型编号的固定顺序完成映射，不使用随机分配，从而保证相同输入始终产生相同的员工级排班结果。

映射完成后，每名员工在十天内始终属于同一小组，每个工作日恰好执行一个完整班型，并恰好工作8天、休息2天。

\subsection{模型求解结果}

由100个日内整数覆盖模型得到各组单日峰值与十天工作人日需求，再根据式\eqref{eq:q1-staff}计算固定人员规模，结果见表\ref{tab:q1-result}。

\begin{table}[H]
\centering
\caption{问题一各小组固定人员规模及最大流结果}
\label{tab:q1-result}
\begin{tabular}{c*{10}{c}}
\toprule
小组 & 1&2&3&4&5&6&7&8&9&10\\
\midrule
单日峰值
&39&39&41&41&40&38&42&41&40&45\\
工作人日下界
&38&38&43&41&44&39&41&42&43&42\\
最优固定编制
&39&39&43&41&44&39&42&42&43&45\\
最大流
&78&78&86&82&88&78&84&84&86&90\\
\bottomrule
\end{tabular}
\end{table}

由表\ref{tab:q1-result}可知，第1、2、7、10组的固定编制由单日峰值决定，第3、5、6、8、9组主要由十天工作人日需求决定，第4组的两项下界相同。十个小组的最优固定编制为
\begin{equation}
\left(
W_1^*,W_2^*,\ldots,W_{10}^*
\right)
=
(39,39,43,41,44,39,42,42,43,45).
\label{eq:q1-result-vector}
\end{equation}

因此，问题一的最少招聘人数为
\begin{equation}
N_1
=
\sum_{g\in G}W_g^*
=
417.
\label{eq:q1-final-result}
\end{equation}

十个小组计算得到的最大流分别为
\[
(78,78,86,82,88,78,84,84,86,90),
\]
均达到对应的$2W_g^*$，说明十个小组均能按照给定的每日工作人数，为每名员工安排两个不同的休息日。

表\ref{tab:q1-roster-sample}给出第1组前两名员工的完整十天排班样例，其中$s_j$表示执行班型$j$，“休”表示当天休息。

\begin{table}[H]
\centering
\caption{问题一员工级排班样例}
\label{tab:q1-roster-sample}
\setlength{\tabcolsep}{5pt}
\begin{tabular}{cccccc}
\toprule
日期 & 第1天 & 第2天 & 第3天 & 第4天 & 第5天\\
\midrule
员工1 & $s_1$ & 休 & $s_1$ & $s_1$ & $s_1$\\
员工2 & 休 & 休 & $s_1$ & $s_1$ & $s_1$\\
\midrule
日期 & 第6天 & 第7天 & 第8天 & 第9天 & 第10天\\
\midrule
员工1 & $s_1$ & $s_1$ & $s_1$ & 休 & $s_1$\\
员工2 & $s_2$ & $s_1$ & $s_1$ & $s_1$ & $s_1$\\
\bottomrule
\end{tabular}
\end{table}

完整排班共包含4170条员工--日期记录，其中3336条为工作记录、834条为休息记录。逐人校验表明，417名员工均恰好工作8天、休息2天，并在十天内固定服务于同一小组；重新计算全部1100个日期--小时--小组需求单元后，缺岗单元数为0，最小覆盖余量为0。完整员工级排班见支撑文件
\texttt{q1\_employee\_schedule.csv}。

作为独立最优性复核，进一步建立包含全部日期、小组和班型变量的完整聚合混合整数规划，SCIP所得最优目标值仍为417，相对最优间隙为0。

\section{问题二模型的建立与求解}

\subsection{统一人员池规模的确定}

问题二允许员工在不同日期服务于不同小组，因此十个固定人员池可以合并为统一人员池。由于员工在同一天内仍只能服务于一个小组，各日期、各小组的日内覆盖模型与问题一相同。

根据式\eqref{eq:q1-daily-cover}得到的日内最低人数，定义第$d$天全场最低工作人数为
\begin{equation}
M_d=\sum_{g\in G}m_{dg}.
\label{eq:q2-daily-lower}
\end{equation}

统一招聘人数一方面不能低于任一日期的全场最低工作人数，另一方面必须能够提供十天所需的全部工作人日。因此
\begin{equation}
W^*
=
\max
\left\{
\max_{d\in D}M_d,
\left\lceil
\frac{\displaystyle\sum_{d\in D}M_d}{8}
\right\rceil
\right\}.
\label{eq:q2-staff}
\end{equation}

计算得到逐日最低工作人数为
\begin{equation}
(M_1,M_2,\ldots,M_{10})
=
(214,349,314,333,349,331,363,350,325,319).
\label{eq:q2-lower-vector}
\end{equation}

其中最高单日需求为363人，十天逐日最低工作人日总和为
\begin{equation}
\sum_{d\in D}M_d=3247.
\label{eq:q2-lower-total}
\end{equation}

由式\eqref{eq:q2-staff}可得
\begin{equation}
W^*
=
\max
\left\{
363,
\left\lceil\frac{3247}{8}\right\rceil
\right\}
=
406.
\label{eq:q2-staff-result}
\end{equation}

因此，问题二的最少招聘人数为406人。

\subsection{整除冗余工作人日的分配}

406名员工每人必须工作8天，因此统一人员池实际需要安排的工作人日为
\begin{equation}
8W^*=406\times8=3248.
\label{eq:q2-required-workdays}
\end{equation}

逐日独立覆盖下界为3247人日，与实际必须安排的工作人日之间相差
\begin{equation}
\Delta
=
8W^*-\sum_{d\in D}M_d
=
1.
\label{eq:q2-redundancy}
\end{equation}

这1个工作人日并非由需求覆盖产生，而是由“每名员工恰好工作8天”的整除条件产生。本文将冗余工作人日优先安排在最低工作人数最大的日期；当多个日期相同时，优先选择日期编号较小者。因此，将第7天的工作人数由363人增加至364人，得到最终每日工作人数
\begin{equation}
(A_1,A_2,\ldots,A_{10})
=
(214,349,314,333,349,331,364,350,325,319).
\label{eq:q2-actual-vector}
\end{equation}

此时
\[
\sum_{d\in D}A_d=3248=8W^*.
\]

在各日期内部，将需要增加的工作人日分配至相应小组，并在固定该小组实际工作人数后，重新求解式\eqref{eq:q1-final-cover}和式\eqref{eq:q1-final-count}，得到各日期、各小组的最终匿名班型配额$x_{dgs}$。

\subsection{员工休息日期的最大流分配}

第$d$天的休息人数为
\begin{equation}
R_d=W^*-A_d.
\label{eq:q2-rest-number}
\end{equation}

由式\eqref{eq:q2-actual-vector}可得
\begin{equation}
\sum_{d\in D}R_d
=
10W^*-\sum_{d\in D}A_d
=
2W^*
=
812.
\label{eq:q2-rest-total}
\end{equation}

建立统一人员池的员工--日期最大流网络。源点到每名员工节点的容量为2，每名员工到每个日期节点的容量为1，第$d$个日期节点到汇点的容量为$R_d$。该网络与问题一中的组内网络结构相同，但员工节点由某一小组的固定员工扩展为全部406名员工。

实际计算得到最大流
\begin{equation}
F=812=2W^*,
\label{eq:q2-full-flow}
\end{equation}
网络达到满流。因此，每名员工均能够分配两个不同的休息日，各日期的在岗人数也与式\eqref{eq:q2-actual-vector}完全一致。

\subsection{员工--小组--班次确定性映射}

对于第$d$天，将全部小组的匿名班型配额展开为
\begin{equation}
\mathcal L_d
=
\bigcup_{g\in G}
\bigcup_{s\in S_0}
\underbrace{\{(g,s),(g,s),\ldots,(g,s)\}}
_{x_{dgs}\text{个}}.
\label{eq:q2-slot-set}
\end{equation}

由班型配额可知
\begin{equation}
|\mathcal L_d|
=
\sum_{g\in G}\sum_{s\in S_0}x_{dgs}
=
A_d.
\label{eq:q2-slot-count}
\end{equation}

最大流确定的第$d$天在岗员工数同样为$A_d$，因此可以将当天在岗员工与$\mathcal L_d$中的匿名槽位进行一一映射。本文按照员工编号、小组编号和班型编号的固定顺序完成配对。

每个匿名槽位只包含一个小组和一个完整班型，因此每名员工在同一天内只会被分配至一个小组，并执行两个属于同一小组的4小时工作段。不同日期之间则允许被分配至不同小组，从而满足问题二的人员调配规则。

\subsection{模型求解结果}

问题二逐日最低工作人数与实际工作人数见表\ref{tab:q2-daily-result}。

\begin{table}[H]
\centering
\caption{问题二逐日最低工作人数与实际工作人数}
\label{tab:q2-daily-result}
\begin{tabular}{c*{10}{c}}
\toprule
日期 &1&2&3&4&5&6&7&8&9&10\\
\midrule
逐日最低人数
&214&349&314&333&349&331&363&350&325&319\\
实际工作人数
&214&349&314&333&349&331&364&350&325&319\\
冗余工作人日
&0&0&0&0&0&0&1&0&0&0\\
\bottomrule
\end{tabular}
\end{table}

表\ref{tab:q2-daily-result}表明，除第7天增加1名工作员工外，其余日期的实际工作人数均等于逐日覆盖下界。该冗余仅用于满足总工作人日的整除要求，不改变最少招聘人数。

表\ref{tab:q2-roster-sample}给出员工56的十天排班样例。“$g/s_j$”表示当天服务于第$g$组并执行班型$j$。

\begin{table}[H]
\centering
\caption{问题二员工级排班样例}
\label{tab:q2-roster-sample}
\setlength{\tabcolsep}{5pt}
\begin{tabular}{cccccc}
\toprule
日期 & 第1天 & 第2天 & 第3天 & 第4天 & 第5天\\
\midrule
员工56
&休&休&$2/s_{10}$&$2/s_4$&$2/s_6$\\
\midrule
日期 & 第6天 & 第7天 & 第8天 & 第9天 & 第10天\\
\midrule
员工56
&$2/s_{10}$&$2/s_4$&$3/s_1$&$2/s_{10}$&$1/s_4$\\
\bottomrule
\end{tabular}
\end{table}

员工56在第1、2天休息，其余8天工作，并分别在第1、2、3组之间进行跨日调配。该样例表明，统一人员池能够在保持员工每天只服务一个小组的条件下，实现不同日期之间的人员共享。

完整员工级排班共包含4060条员工--日期记录，其中3248条为工作记录、812条为休息记录。校验结果表明：

\begin{enumerate}[label=\textbf{(\arabic*)},leftmargin=4em]
  \item 406名员工均恰好工作8天、休息2天；
  \item 每名员工在一个工作日内只服务于一个小组；
  \item 各日期实际工作人数与式\eqref{eq:q2-actual-vector}完全一致；
  \item 各日期、各小组的实际班型人数与匿名班型配额逐项一致；
  \item 1100个需求单元均不存在人员缺口，最小覆盖余量为0。
\end{enumerate}

完整排班见支撑文件
\texttt{q2\_employee\_schedule.csv}。

\subsection{独立复核}

为检验解析招聘规模与最大流构造结果，本文分别采用完整聚合混合整数规划和CP-SAT进行独立复核。

完整聚合模型直接以招聘人数为目标，同时包含全部日期、小组、班型覆盖及工作人日守恒约束。SCIP求得问题一和问题二的最优目标值分别为417和406，相对最优间隙均为0，与分解算法所得结果完全一致。

进一步设$y_{id}=1$表示员工$i$在第$d$天工作，否则为0，并建立
\begin{align}
\sum_{d\in D}y_{id}&=8,
&&\forall i,\label{eq:cpsat-employee}\\
\sum_i y_{id}&=A_d,
&&\forall d.\label{eq:cpsat-day}
\end{align}

采用CP-SAT对上述工作--休息矩阵进行独立可行性检验
\cite{googlecpsat}。问题一的十个固定人员池均得到可行解，问题二同样得到可行解，且逐日工作人数与最大流主构造完全一致。因此，解析招聘规模、最大流休息日分配和员工级班次映射之间保持一致。

\subsection{两种人员组织方式的比较}

问题一和问题二的最少招聘人数分别为417人和406人。允许员工在不同日期之间更换服务小组后，招聘人数减少
\begin{equation}
N_1-N_2=417-406=11,
\end{equation}
相对降幅为
\begin{equation}
\frac{417-406}{417}\times100\%
=
2.64\%.
\end{equation}

两种方案均未改变逐小时需求和每名员工每天工作8小时、十天内做8休2的规定。招聘人数的减少来自不同小组需求峰值在日期上的错位：固定归组时，各小组必须分别按照自身峰值和工作人日需求配置人员；合并为统一人员池后，员工可以在不同日期流向需求较高的小组，从而降低固定人员池中的闲置。
\section{问题三模型的建立与求解}

\subsection{结构性死锁与规则修复}

\subsubsection{公共盲区定理}

设营业开始、结束时刻为$O,C$，窗口长度$L=C-O$；两段工作长度均为$b$，最小休息为$r$。第一段为了给第二段留出$b+r$小时，最晚结束于$C-b-r$；第二段最早开始于$O+b+r$。若
\begin{equation}
2(b+r)>L,
\label{eq:blindcondition}
\end{equation}
则两者之间存在所有合法班型均无法覆盖的公共盲区，其长度为
\begin{equation}
B=2(b+r)-L.
\label{eq:blindlength}
\end{equation}

本题$O=8,C=19,L=11,b=4,r=2$，故$B=1$小时。严格合法班型只有：
\begin{equation}
\begin{aligned}
&(08{:}00{-}12{:}00,\ 14{:}00{-}18{:}00),\\
&(08{:}00{-}12{:}00,\ 15{:}00{-}19{:}00),\\
&(09{:}00{-}13{:}00,\ 15{:}00{-}19{:}00).
\end{aligned}
\end{equation}
如图\ref{fig:blind}所示，三种班型在13:00--14:00的覆盖均为0。

\begin{figure}[H]
  \centering
  \originalfigure[width=0.94\textwidth]{fig_strict_blind_zone.png}
  \caption{严格休息规则下全部合法班型及公共盲区}
  \label{fig:blind}
\end{figure}

附件在该小时每天十组总需求依次为85、98、146、149、152、159、188、206、150、167人，十天合计1500人次。设$z_{dsgk}$为前段服务组$g$、后段服务组$k$的全职人数，则盲区覆盖约束左端恒为0，而右端$q_{d,5,g}\ge1$，产生$0\ge q_{d,5,g}$的矛盾。因此严格问题三不是“当前人数不够”，而是\textbf{任意有限招聘人数均不可行}。

\subsubsection{修复方案一：将最小休息缩短为1小时}

把班型集合由$S_2$放宽为$S_1$，合法班型从3种增至6种，13:00--14:00不再是共同盲区。跨组覆盖约束写为
\begin{equation}
\sum_{s\in S_1}\left[
\ind\{u_s\le h<u_s+4\}\sum_{k\in G}z_{dsgk}
+\ind\{v_s\le h<v_s+4\}\sum_{j\in G}z_{dsjg}
\right]\ge q_{dhg}.
\label{eq:q3crosscover}
\end{equation}
保持全职做8休2，整数最优为437人。该方案数学上最直接，但相较无最小间隔对照的400人增加37人，而且把至少2小时休息压缩为1小时，员工疲劳代价没有被显式计入人数目标。

\subsubsection{修复方案二：允许非对称工作段}

进一步定义非对称班型$s=(u,L_1,v,L_2)$，满足
\begin{equation}
L_1,L_2\le4,\qquad v\ge u+L_1+2,qquad v+L_2\le11.
\label{eq:asymshift}
\end{equation}
为避免过度缩短，数值比较仅允许一个工作段由4小时缩为3小时，形成4+3和3+4两类；2小时工作段作为可扩展边界保留，不在主结果中启用。采用字典序目标：先最小化非4+4班型的员工人日，再固定该最小值并最小化招聘人数。只允许4+3时需495人；同时允许3+4时降至439人。两种方案均需1500个修复员工人日，恰对应盲区需求下界，但员工实际工作7小时却按全天编制计费，造成带薪闲置。

\subsubsection{修复方案三：严格全职与午间兼职的混合用工}

主推方案保留全职员工的$S_2$严格班型，引入覆盖13:00--14:00的4小时兼职班。允许兼职起点$t\in T=\{10,11,12,13\}$，设$a^P_{ht}$表示该兼职班是否覆盖小时$h$，$p_{dtg}$为其人数。联合覆盖为
\begin{equation}
\sum_{s\in S_2}\left[
\ind\{u_s\le h<u_s+4\}\sum_{k}z_{dsgk}
+\ind\{v_s\le h<v_s+4\}\sum_{j}z_{dsjg}
\right]
+\sum_{t\in T}a^P_{ht}p_{dtg}\ge q_{dhg}.
\label{eq:mixedcover}
\end{equation}
全职人员池满足
\begin{align}
\sum_{s,g,k}z_{dsgk}+R_d&=W^F, &&\forall d,\label{eq:mixedday}\\
\sum_dR_d&=2W^F.\label{eq:mixedrest}
\end{align}
定义十天兼职班总人次
\begin{equation}
P=\sum_{d\in D}\sum_{t\in T}\sum_{g\in G}p_{dtg}.
\end{equation}
由于严格全职在13:00--14:00的覆盖为0，每个兼职班在该小时至多贡献1人，故
\begin{equation}
P\ge\sum_{d\in D}\sum_{g\in G}q_{d,5,g}=1500.
\label{eq:partlower}
\end{equation}

为准确表达“兼职只修补缺口，而非替代全部全职”，采用字典序目标：
\begin{equation}
\text{第一层：}\min P;\qquad
\text{第二层：在 }P=P^*\text{ 下 }\min W^F.
\label{eq:lexico}
\end{equation}
第一层得到$P^*=1500$，达到式\eqref{eq:partlower}的理论下界；第二层得到$W^F=325$。这证明325不是预设参数，而是“最小兼职规模”政策下的全局最优全职编制。每日全职与兼职配置见图\ref{fig:mixedday}。

\begin{figure}[H]
  \centering
  \originalfigure[width=0.92\textwidth]{fig_mixed_daily.png}
  \caption{混合用工方案的每日全职在岗人数与兼职班人次}
  \label{fig:mixedday}
\end{figure}

需要说明，325是\textbf{全职招聘人数}，1500是\textbf{十天累计兼职班人次}，二者不能直接相加为1825名唯一人员。若兼职可跨日重复上岗，则唯一兼职人数取决于后续合同设计；仅按单日峰值，至少需要容纳第8天的206个兼职班。

\subsubsection{三类修复的二维评价}

在同小时工资口径下，一名全职会期成本折算为64小时，一个兼职班折算为4小时，混合方案等效全职成本为
\begin{equation}
C_{\mathrm{eq}}=W^F+\frac{P}{16}=325+\frac{1500}{16}=418.75.
\label{eq:eqcost}
\end{equation}
定义付费工时浪费率
\begin{equation}
\eta=\frac{H_{\mathrm{paid}}-17682}{H_{\mathrm{paid}}}.
\label{eq:waste}
\end{equation}
各方案比较见表\ref{tab:repair}和图\ref{fig:repairmatrix}。非对称班型按全天工资计费，因此缩短的工作小时也视为付费闲置。

\begin{table}[H]
\centering
\caption{问题三规则修复方案比较}
\label{tab:repair}
\begin{tabularx}{\textwidth}{L>{\centering\arraybackslash}p{2.3cm}>{\centering\arraybackslash}p{2.3cm}>{\centering\arraybackslash}p{2.2cm}L}
\toprule
方案 & 招聘或等效成本 & 付费工时 & 浪费率 & 规则代价\\
\midrule
休息缩短为1小时 & 437人 & 27968 & 36.78\% & 疲劳风险增加\\
仅允许4+3 & 495人 & 31680 & 44.19\% & 大量带薪早退\\
允许4+3与3+4 & 439人 & 28096 & 37.07\% & 带薪早退\\
全职+午间兼职 & 418.75等效人 & 26800 & 34.02\% & 需管理兼职池\\
无最小间隔（对照） & 400人 & 25600 & 30.93\% & 不满足严格休息\\
\bottomrule
\end{tabularx}
\end{table}

\begin{figure}[H]
  \centering
  \originalfigure[width=0.86\textwidth]{fig_repair_cost_waste.png}
  \caption{规则修复方案的成本--浪费率二维评价}
  \label{fig:repairmatrix}
\end{figure}

混合方案位于严格可行方案中的左下区域：相较437人的休息1小时方案，等效成本下降18.25人，浪费率下降2.76个百分点；相较439人的双向非对称方案，等效成本下降20.25人，且不缩短全职员工的每日实际工作要求。无最小间隔对照虽然数值最低，但违反问题三的休息制度，不能作为正式答案。

\section{模型检验与灵敏度分析}

\pending{待补充：求解器版本、运行环境、硬件配置、每个模型的运行时间及完整求解日志。当前临时稿仅给出求解状态和结果复核。}

\subsection{输入、最优状态与覆盖复核}

数据预检确认附件包含110条有效日期--小时记录、1100个正需求单元，输入文件SHA-256固定为\texttt{9c65acae...98d44}。问题一、问题二、休息1小时修复、无最小间隔对照、非对称班型与混合用工均由SCIP返回OPTIMAL；严格问题三同时由解析盲区与整数模型判为不可行。所有输出在运行后记录脚本、输入和图表哈希，以防正文数字与代码结果脱节。

对于宏观解，重新计算每个$d,h,g$的覆盖余量
\begin{equation}
e_{dhg}=\text{实际覆盖}_{dhg}-q_{dhg}.
\end{equation}
全部$e_{dhg}\ge0$。问题二员工级排班的最小余量为0，说明存在恰好贴住需求的单元，但没有缺岗。

\subsection{宏观--微观一致性检验}

问题二包含三层对账：第一，宏观每日工作人数之和为3248，等于$8\times406$；第二，CP-SAT生成的每名员工行和均为8、休息数均为2；第三，贪心映射后每一天的实际班次数与$x_{dgs}$逐项相等。上述对账排除了“总人日正确但个人工作天数错误”和“状态正确但班次覆盖错误”两类常见漏洞。

\subsection{休息口径敏感性}

为检验417和406对语义的依赖，在问题一、二中额外强加至少1小时休息并重新求解，人数分别上升至524和484。这一结果不作为题目答案，只说明最小休息口径会显著收缩班型集合。正式模型因此遵循原题：问题一、二使用$S_0$，问题三使用$S_2$，休息1小时仅作为问题三修复方案。

\subsection{兼职工资权重敏感性}

若不采用“先控制兼职规模”的字典序政策，而直接最小化$16W^F+\lambda P$，全职与兼职会随$\lambda$发生替代。图\ref{fig:wagesens}显示，当兼职成本权重从1增加到16时，全职最优规模由283升至321，兼职班人次由1864降至1504。该结果说明325并非所有工资结构下的成本最优点，而是\textbf{最少兼职优先}政策下的最优点；论文将政策条件与数值结论绑定，避免脱离管理偏好宣称唯一最优。

\begin{figure}[H]
  \centering
  \originalfigure[width=0.86\textwidth]{fig_wage_sensitivity.png}
  \caption{兼职工资权重变化下的全职--兼职替代关系}
  \label{fig:wagesens}
\end{figure}

\section{模型评价与推广}

\subsection{模型优点}

\textbf{第一，结构解释强。} 对问题三给出公共盲区定理，证明无解与招聘人数无关，优于只报告求解器不可行。\textbf{第二，计算与执行分层。} 宏观MILP专注最少人数，微观CP-SAT专注个人规则，既降低维度又保留员工级可执行性。\textbf{第三，管理含义清楚。} 417到406量化跨天共享价值；437、495、439与混合方案揭示不同制度修复的成本。\textbf{第四，结果可审计。} 题目数据、脚本、结果、图表和完整员工表均可通过运行记录追溯，而不是由文字推测数字。

\subsection{模型不足}

本文把小时需求视为确定值，未考虑客流预测误差与临时缺勤；员工能力同质，未区分技能等级、培训成本和组别偏好；兼职以人次计量，尚未决定同一兼职员工可连续承担多少天；等效成本采用同小时工资，未加入最低班次补贴、招聘固定成本、夜班溢价与管理成本。贪心映射保证硬约束，但未进一步优化连续晚班、跨组频率和员工偏好。

\subsection{改进方向}

可在现有双层框架上扩展：一是在宏观层加入多情景需求和预算不确定集，形成稳健招聘模型；二是在微观层固定最少人数后，引入晚班均衡、连续工作天数和换组次数的字典序目标；三是为兼职建立跨日身份变量，将1500人次转化为唯一兼职人数和合同组合；四是对不同小组建立技能兼容矩阵，以最少交叉培训边获得接近统一人员池的节约。

\subsection{推广应用}

本文方法适用于会展、景区、呼叫中心、零售门店、医院辅助岗位和大型活动安保等“需求按小时变化、人员跨日共享、班型受休息规则限制”的场景。员工排班研究表明，现实系统还需同时兼顾服务水平、劳动协议、偏好和公平性\cite{ernst2004}。双层框架允许宏观层更换成本和需求模型，微观层增加个人约束，而不必重写整个模型。

\section{结论}

本文围绕展销会三种工作模式建立了一条统一主线：先枚举合法班型并求匿名班次数量，再把匿名解落实到具体员工。问题一的固定小组人员池需417人；问题二借助跨天共享降至406人，并成功构造每人做8休2的完整员工级排班。问题三在严格4小时工作、至少2小时休息、再工作4小时的规则下存在13:00--14:00公共盲区，因此原问题结构性无解。

在规则修复中，休息缩短为1小时需437人，非对称班型分别需495或439人。主推混合用工方案保持全职待遇，以午间4小时兼职班填补盲区；字典序优化首先把兼职压到理论下界1500人次，再将全职编制压到325人，等效成本为418.75名全职、浪费率为34.02\%。这一结果表明，面对严格劳动规则造成的不可行性，最有效的策略不是无差别增加全职，而是识别结构性缺口并引入边界明确的局部柔性。

\bibliographystyle{gbt7714-numerical}
\bibliography{references}

\appendix
\section{支撑材料与可复现说明}

\pending{当前仅完成论文正文迁移。原始附件、完整程序、结果 JSON/CSV、求解日志和原始图片尚未随本次临时稿提供，正式提交前必须补齐。}

\subsection{主要文件}

\begin{longtable}{p{0.35\textwidth}p{0.56\textwidth}}
\toprule
文件 & 用途\\
\midrule
\endfirsthead
\toprule
文件 & 用途\\
\midrule
\endhead
\path{solve_core_and_roster.py} & 【待补充】问题一、问题二、严格问题三、休息修复及员工级重建\\
\path{solve_asymmetric_repair.py} & 【待补充】4+3与3+4非对称班型字典序优化\\
\path{solve_mixed_workforce.py} & 【待补充】严格全职与午间兼职的联合模型、325人结论与工资敏感性\\
\path{run_modeling.py} & 【待补充】统一运行全部模型与绘图\\
\path{q2_employee_schedule.csv} & 【待补充】406名员工十天完整排班，共4060条记录\\
\path{core_model_results.json} & 【待补充】417、406、437、400与员工级校验结果\\
\path{mixed_workforce_results.json} & 【待补充】325名全职、1500兼职人次及敏感性结果\\
\bottomrule
\end{longtable}

\subsection{公共盲区证明的边界}

定理要求两段工作长度相同且最小间隔固定。当$2(b+r)\le L$时，式\eqref{eq:blindlength}不产生正盲区，但这只说明不存在由窗口长度必然造成的公共空档，不自动保证给定需求一定可覆盖。若不同员工允许不同段长，需重新取所有合法班型覆盖集合的交集。

\subsection{结果复现实验顺序}

\begin{enumerate}[label=\textbf{Step \arabic*:},leftmargin=6em]
  \item 保持原始附件只读，确认文件哈希；
  \item 运行统一入口生成核心、非对称和混合用工结果；
  \item 检查求解状态均为OPTIMAL或解析不可行；
  \item 校验406名员工工作日、休息日和逐小时覆盖；
  \item 由结果JSON重新绘图并生成方案比较表；
  \item 编译论文并逐页检查公式、图表和引用。
\end{enumerate}



\end{document}
